\chapter{Conclusions}
\label{c:conclusions}
\section{Conclusió genèrica}
Aquest treball de recerca sorgeix d'una necessitat concreta que tenen alguns alumnes que han cursat Matemàtiques II (MAT) i que, posteriorment, es volen presentar a les PAU de Matemàtiques Aplicades a les Ciències Socials II (MACS), com jo mateix. Per això vaig decidir investigar quins coneixements calia adquirir i, a partir d'aquí, elaborar una guia que facilites preparar-se per l'examen tal i com s'ha explicat al capítol~\ref{recerca previa} \nameref{recerca previa}.

Després de fer la recerca, estudiar els continguts i resoldre exercicis i exàmens de les PAU, he pogut acabar la guia que em vaig proposar fer per estudiants que es trobin en aquesta situació. L'objectiu del treball s'ha assolit, ja que la guia permet identificar quins continguts de MACS no s'han treballat a MAT i proporciona els recursos teòrics i práctics per a poder afrontar la prova de MACS.

La guia ha estat dissenyada per a tots els alumnes de MAT que es vulguin presentar a la prova de MACS. A més de recomanar-la per a aquests estudiants, també pot resultar util per al alumnes de MACS que volen complementar i/o consolidar la part de probabilitat i estadística.

Elaborar aquesta guia m'ha permès ampliar els meus coneixements sobre les matemàtiques. De fet, vaig fer dos exàmens d'anys anteriors i els he fet sense cap problema. Aquest resultat em fa pensar que estudiar els continguts que em faltaven i practicar amb exàmens reals m'ha servit, no només per elaborar la guia, sinó també per assolir un altre objectiu: poder afrontar amb un examen de MACS de les PAU.

En conclusió, la finalitat i els objectius marcats per a aquest treball s'han assolit sobradament. He pogut estudiar continguts que necessitava, posar a prova els coneixements, i elaborar una guia transmeten els coneixements perquè altres alumnes els hi siguin d'utilitat.

\section{Tancament personal}
Aquest projecte ha estat el meu primer treball de recerca. A Catalunya, habitualment, durant quart d'ESO se'n fa un, però jo no el vaig fer perquè vaig cursar quart d'ESO a Irlanda. Per això, afrontar aquest treball ha estat una nova experiència en aquest tipus de projectes.

Una de les coses que més m'ha agradat ha estat poder decidir com enfocar el treball i organitzar-me pel meu compte. Al principi, però, això em costava: moltes vegades em quedava encallat i no sabia quin havia de ser el pas següent. El meu tutor sempre va estar disponible per a qualsevol cosa que li demanés, però com que no el veia cada dia, i a més preferia intentar solucionar els problemes jo mateix abans de preguntar-li, sovint tirava endavant sol.

Quan em vaig marcar els objectius volia aprofundir en les matemàtiques. Mentrestant, estudiava els temes i els explicava, m'ho passava bé. No m'he limitat a estudiar només els temes que entren a les PAU, també m'he mirat temes com la Paradoxa de Sant Petersburg plantejada per Nicolaus Bernoulli el 1713 o la teoria dels jocs (que combina les matemàtiques amb l'estudi de la presa de decisions) de von Neumann.

M'agrada que el treball tingui una utilitat concreta i ajudi a molta gent per la preparació per les PAU de MACS. Un cop acabada la guia i veient tot el treball fet, em sento molt realitzat. D'altra banda, també m'ha resultat útil treballar amb temes didàctics, perquè soc professor d'escacs de nens petits i haver pogut fer una guia didàctica m'ha ajudat a ampliar els meus coneixements del camp, sobretot, en millorar les planificacions amb més claredat dels temes que es volen explicar.

\section{Recerca futura}
Si tingués més temps per a continuar fent el treball, m'agradaria aprofundir en la branca de probabilitat i orientar una futura recerca cap als jocs d'atzar. De fet, abans de decidir el tema definitiu del treball de recerca, un dels altres candidats que vaig considerar va ser investigar algun aspecte matemàtic relacionat amb els jocs d'atzar. Aquest àmbit em crida especialment l'atenció perquè permet aplicar conceptes matemàtics a situacions que podem trobar en la vida quotidiana i, a la vegada, planteja preguntes sobre el paper de l'atzar i la presa de decisions, posant una mica de psicologia al joc i barrejant-la amb les matemàtiques.

Una possible línia de recerca seria analitzar matemàticament diferents jocs d'atzar, com el blackjack, el pòquer o altres jocs de casino, i estudiar en quines condicions un jugador té avantatge o desavantatge. En lloc de quedar-me només amb les regles, m’agradaria calcular les probabilitats que hi ha en cada cas, veure quin és el valor esperat de les apostes i comparar-ho amb l’avantatge que té estadísticament la banca. Això permetria investigar fins a quin punt les decisions d'un jugador poden influir en el resultat i quina part correspon realment a l'atzar.

A partir d’aquesta anàlisi també es podria estudiar una qüestió que em sembla interessant: fins a quin punt un jugador pot obtenir resultats favorables repetidament en un joc basat en ``l’atzar''.

Per tant, una futura recerca podria centrar-se a estudiar matemàticament els jocs d’atzar, amb la probabilitat i l’estadística per entendre millor el funcionament, i inclús, fer alguna guia sobre quan un jugador té més avantatge o sobre com augmentar les seves possibilitats de guanyar en aquest tipus de jocs. Crec que podria ser una continuació del treball, ja que em permetria aprofundir en una part de les matemàtiques, la probabilitat, en què he treballat durant l'elaboració d'aquest projecte i, al mateix temps, investigar algunes de les qüestions que ja em plantejava abans de començar el TR.
